\section{Implementation}
\label{sec:implementation}

\tool has three parts: a tracer and a replay library in C, both interposed through MPI's
profiling interface, and an analyser in Python.

\paragraph{Tracer.} The tracer is preloaded into the unchanged binary. It writes one JSON
record per call and rank: the arguments, the result, the request it creates or completes,
and the communicator as a global identifier with its members' world ranks. Blocking calls
first emit a \emph{begin} record, so that a rank killed while blocked shows where it was
blocked. Repeated failed polls (\mpicall{Iprobe}, \mpicall{Test}) with no call in between
are written once with a count. Each record from user code names its call site, the return
address of the MPI call as an object and an offset, which are the same in every process; the
report groups the findings by the site of the changed receive or probe, and \texttt{addr2line}
names a site in the source. Derived datatypes are flattened to their basic type
signature for the count and type checks of \cls{1a}. The tracer covers all point-to-point
modes, persistent requests, \mpicall{Sendrecv}, every wait and test variant, probes
including matched probes, cancellation, blocking and nonblocking collectives, and
communicator and intercommunicator creation. Under \texttt{MPI\_THREAD\_MULTIPLE},
records carry a thread identifier. The analysis then reports observed findings only,
because a rank's program order is no longer total. Calls outside this fragment (RMA,
neighbourhood collectives, dynamic processes, generalised requests, and the MPI-4
\mpicall{Isendrecv}, partitioned and large-count calls) write an \emph{unmodelled} record.
The analysis then reports observed findings only, because the trace is no longer a complete
record of the run's communication. \pending{C-2: record results, and side records written
after completions.}

\paragraph{Analyser.} The analyser rebuilds $\Mobs$ from the statuses and checks the
recorded run itself (\cref{tab:classes}, the observed classes). It then enumerates
steals and swaps by queue heads. A steal at a blocking receive is decided by a fast check.
The horizon is computed by a worklist over ranks with suffix-minimum indexes per pair of
ranks. Rule~\ref{H:changed} cuts the rank right after the receive, so the new match has one
possible kept successor, a rendezvous completion. Acyclicity therefore reduces to a few
reachability queries in the graph of the recorded run, answered by vector clocks computed
once per threshold. Steals at nonblocking receives and swaps go to a generic check, which
builds the horizon and the graph from scratch. The generic check is the reference, and the
tests compare the two on every alternative of every traced idiom. The \LoneSP explorer runs below a size limit. Two-change
alternatives on different ranks are available on request.
\pending{Cost per alternative after T4; tree clocks~\citep{mathur22treeclock} as a
drop-in for the vector clocks.}

\paragraph{Replay.} The replay library reads the pins from the environment and
substitutes the source of the pinned wildcard receives and probes, persistent receives
included, through a shadow request. It can also delay chosen sends and receives, which the
stress-testing baseline of \cref{sec:eval} uses. When a pinned receive completes or a pinned
probe finds a message, the library logs the match or the envelope (\cref{sec:replay}).
